%TCIDATA{Version=5.50.0.2890}
%TCIDATA{LaTeXparent=0,0,NoticeDKE.tex}


We recall the electron drift kinetic equation may be expressed as
\begin{equation}
v_{s}\frac{\partial f}{\partial s}+\frac{v_{\Vert }}{\Omega }I\left( \psi
\right) \frac{\left| \nabla \psi \right| }{R}\frac{\partial }{\partial s}%
\left( \frac{v_{\Vert }}{B}\right) \frac{\partial f}{\partial \psi }=%
\mathcal{C}\left( f\right) +\mathcal{Q}\left( f\right) +\mathcal{E}\left(
f\right)
\end{equation}

Each of these terms corresponds to a time evolution, and is therefore
associated with a time-scale:

\begin{itemize}
\item  \textit{Motion along magnetic field lines}. The time scale here is
the bounce time $\tau _{b}$ for trapped electrons, or the transit time $\tau
_{t}$ for circulating ones, which can be deduced directly from expression (%
\ref{Normbtime}). For circulating electrons
\begin{equation}
\tau _{t}=\frac{2\pi R_{p}\widetilde{q}}{v\left| \xi _{0}\right| }\sim \frac{%
2\pi qR_{p}}{v_{Te}}
\end{equation}

taking $\lambda \sim 1$ in that case, and  $q\simeq \widetilde{q}$ that is
valid circular plasma cross-sections, and $v\left| \xi _{0}\right| =v_{\Vert
}\simeq v_{Te}$ for thermal electrons. Since $v_{\Vert }\leq \sqrt{\epsilon }%
v_{\bot }$ $\simeq \sqrt{\epsilon }v_{Te}$ for trapped electrons, as the
consequence of the magnetic moment conservation, it turns out that $\tau _{b}
$ may be deduced directly from $\tau _{t}$%
\begin{equation}
\tau _{b}=\frac{2\pi R_{p}\widetilde{q}}{v\left| \xi _{0}\right| }\sim \frac{%
2\pi qR_{p}}{\sqrt{\epsilon }v_{Te}}
\end{equation}

Consequently, $\tau _{b}\geq \tau _{t},$ since $\epsilon \leq 1,$ a result
which is the consequence of the progressive slowing down of the parallel
velocity as far as the electron approaches the turning point.

\item  \textit{Vertical drift across magnetic field lines}. The time scale
here is the drift time $\tau _{d}$, which corresponds to the time for an
electron to drift across the plasma. It is then given by
\begin{equation}
\tau _{d}=\int_{0}^{\psi _{a}}d\psi \frac{\Omega _{e}}{v_{\Vert }I\left(
\psi \right) }\frac{R}{\left| \nabla \psi \right| }\left[ \frac{\partial }{%
\partial s}\left( \frac{v_{\Vert }}{B}\right) \right] ^{-1}
\end{equation}
using expression \ref{drift3} for the radial component $\mathbf{v}_{D}\cdot %
\nabla \psi $ of the drift velocity$\mathbf{\ v}_{D}$ as defined in \ref
{vdrift}. Using relation  $I\left( \psi \right) $ as defined in (\ref{Bt1}),
and the fact that $d\psi /\left| \nabla \psi \right| =dr$ for circular
concentric plasma cross-sections,
\begin{equation}
\tau _{d}\simeq \frac{\Omega _{e}}{v_{Te}^{2}}\int_{0}^{a_{p}}\frac{dr}{B_{T}%
}\left[ \frac{\partial }{\partial s}\left( \frac{1}{B}\right) \right] ^{-1}
\end{equation}
and since in that magnetic configuration $\partial /\partial s\sim 2\pi %
/R_{p},$ ones obtains finally
\begin{equation}
\tau _{d}\simeq 2\pi \frac{\Omega _{e}}{v_{Te}^{2}}R_{p}\int_{0}^{a_{p}}%
\frac{B}{B_{T}}dr
\end{equation}
or
\begin{equation}
\tau _{d}\simeq 2\pi \frac{\Omega _{e}}{v_{Te}^{2}}R_{p}a_{p}
\end{equation}
Consequently, the small drift parameter $\delta _{d}$ is defined as the ratio
\begin{equation}
\delta _{d}\equiv \frac{\tau _{t}}{\tau _{d}}\sim q\frac{\rho _{L}}{a_{p}}%
\sim \frac{\rho _{L}}{a_{p}}  \label{smalldriftordering}
\end{equation}
where the thermal Larmor radius $\rho _{L}=$ $v_{Te}/\Omega _{e}$ has been
introduced.

\item  \textit{Collisions}. The Fokker-Planck theory considers the
cumulative effect of many small-angle collisions in calculating the rate of
change of a particle distribution. From this theory, the characteristic time
scale $\tau _{c}$ corresponds for Coulomb collisons to deflect an electron's
path by a significant angle, on the order of $\pi /2.$ This is the electron
thermal collision time $\tau _{c}$ whose expression is
\begin{equation}
\tau _{c}=\nu _{c}^{-1}=\frac{4\pi \varepsilon _{0}m_{e}^{2}v_{Te}^{3}}{%
q_{e}^{4}n_{e}\ln \Lambda }
\end{equation}
where $\ln \Lambda $ is the well known Coulomb logarithm, a slowly varying
function of the plasma temperature and density. The collision time scale $%
\tau _{c}$ holds for circulating electrons.

For trapped ones, it is more physical to consider an alternate collision
time scale determined not by the time for deflection of the path by $\pi /2,$
but by the time needed for the electron to be deflected so that it is no
longer on a trapped orbit. In the limit $\epsilon \ll 1,$ we can approximate
the change of the pitch angle necessary to make trapped particles become
untrapped
\begin{equation}
\Delta \xi _{0}\sim \xi _{0T}\simeq \sqrt{\frac{2\epsilon }{1+\epsilon }}%
\sim \sqrt{\epsilon }
\end{equation}
Because the small-angle collisions produce a random-walk change in the
pitch-angle $\xi _{0},$ the \textit{effective} collision time for detrapping
$\tau _{c}^{eff.}$ or $\tau _{dt}$ is approximately deduced from relation
\begin{equation}
\nu _{c}^{eff.}\sim \frac{\nu _{c}}{\Delta \xi _{0}^{2}}\sim \frac{\nu _{c}}{%
\epsilon }
\end{equation}
and
\begin{equation}
\tau _{dt}=\tau _{c}^{eff.}\sim \epsilon \tau _{c}
\end{equation}

\item  \textit{Constant electric field acceleration}. From the relation (\ref
{EfieldFPterm}), it is straightforward to estimate that the time scale
associated to constant electric field acceleration is
\begin{equation}
\tau _{e}\sim \frac{m_{e}v_{Te}}{\left| q_{e}\right| E_{\Vert }}\sim \tau
_{c}\frac{\left| E_{D}\right| }{E_{\Vert }}
\end{equation}
where the well known Dreicer field $E_{D}$%
\begin{equation}
E_{D}=\nu _{c}\frac{m_{e}v_{Te}}{q_{e}}
\end{equation}
is introduced. Here only circulating particle are concerned.

\item  \textit{Quasilinear diffusion}. In a similar approach,
\begin{equation}
\tau _{ql}\sim \frac{\left( m_{e}v_{Te}\right) ^{2}}{D_{ql}}\sim \tau _{c}%
\frac{D_{ql}^{\dagger }}{D_{ql}}
\end{equation}

where $D_{ql}^{\dagger }=\nu _{c}\left( m_{e}v_{Te}\right) ^{2}.$
\end{itemize}

\subsection{Small Drift Ordering}

In the small drift expansion $\delta _{d}\ll 1,$ where the small parameter
is defined by relation (\ref{smalldriftordering})
\begin{equation}
f=f_{0}+f_{1}+\cdots
\end{equation}
the first order equation is
\begin{equation}
v_{s}\frac{\partial f_{0}}{\partial s}=\mathcal{C}\left( f_{0}\right) +%
\mathcal{Q}\left( f_{0}\right) +\mathcal{E}\left( f_{0}\right)
\end{equation}
which is usually referred to as the Fokker-Planck equation.

The second order equation is
\begin{equation}
v_{s}\frac{\partial f_{1}}{\partial s}+\frac{v_{\Vert }}{\Omega _{e}}I\left(
\psi \right) \frac{\left| \nabla \psi \right| }{R}\frac{\partial }{\partial s%
}\left( \frac{v_{\Vert }}{B}\right) \frac{\partial f_{0}}{\partial \psi }=%
\mathcal{C}\left( f_{1}\right) +\mathcal{Q}\left( f_{1}\right) +\mathcal{E}%
\left( f_{1}\right)
\end{equation}
which is referred to as the electron Drift Kinetic equation.
